\documentclass[preprint,12pt,a4paper]{elsarticle}
\usepackage{amsmath,amssymb,booktabs,tabularx,hyperref,lineno}
\usepackage{microtype}
\setlength{\parindent}{0pt}
\setlength{\emergencystretch}{3em}
% Public source revision; availability must be verified separately.
\newcommand{\artifactroot}{https://github.com/imrangadzhiev/reaction-lens-systems/blob/14c0ac048777396b121ea877df4f1a08df935a66}
\newcommand{\repositorynotice}{Code and data: \url{https://github.com/imrangadzhiev/reaction-lens-systems}.}
\newcommand{\evidence}[2]{\href{\artifactroot/#1}{#2}}

\input{numbers.tex}
\journal{SoftwareX}
\begin{document}
\begin{frontmatter}
\title{Reaction Lens Systems: Reproducible Full-Catalog Retrieval of Reactome Reactions from Scientific Text}
\author{Imran Gadzhiev\corref{cor1}}
\ead{i.gadzhiev.m@gmail.com}
\cortext[cor1]{Corresponding author.}
\address{Independent Researcher, Manchester, United Kingdom}
\begin{abstract}
Reaction Lens Systems is open-source research software that ranks Reactome
reactions from English scientific titles and abstracts. A frozen scientific-text
encoder and learned global/local scoring head evaluate all 15,692 reactions. The release
provides pinned model artifacts, candidate-transformation checks, component
timings and reproducible numerical evidence. In recorded CPU experiments,
candidate transformations and blocking preserve sampled pairwise scores, while
text encoding dominates complete prediction cost. The software supports
inspectable retrieval proposals for human review and verification of released-model
ranking metrics against recorded Reactome links.
\end{abstract}
\begin{keyword}
Scientific retrieval \sep Reactome \sep Research software \sep Reproducibility \sep Text mining
\end{keyword}
\end{frontmatter}
\linenumbers
\section*{Required Metadata}
\section*{Current code version}
\begin{table}[ht]
\small
\begin{tabularx}{\textwidth}{|l|p{0.35\textwidth}|X|}
\hline
\textbf{Nr.} & \textbf{Code metadata description} & \textbf{Value} \\
\hline
C1 & Current code version & 0.1.0.dev0; measured source revision \texttt{2140cc40ba3d} (full identifier in C2) \\
\hline
C2 & Permanent GitHub link to code/repository used for this code version & \href{https://github.com/imrangadzhiev/reaction-lens-systems/tree/2140cc40ba3dc2af9ad0600cd2cc1624a47a2bf8}{GitHub: pinned source revision} \\
\hline
C3 & Legal Code License & MIT; third-party model and data terms documented separately \\
\hline
C4 & Code versioning system used & git \\
\hline
C5 & Software code languages, tools, and services used & Python, PyTorch, NumPy; FastAPI/Uvicorn local service; Hugging Face artifact hosting \\
\hline
C6 & Compilation requirements, operating environments \& dependencies & Python 3.11--3.12; dependencies pinned in uv.lock; evaluated on macOS arm64 CPU FP32 \\
\hline
C7 & If available Link to developer documentation/manual & \href{https://github.com/imrangadzhiev/reaction-lens-systems/blob/2140cc40ba3dc2af9ad0600cd2cc1624a47a2bf8/README.md}{README: installation, usage and verification} \\
\hline
C8 & Support email for questions & \href{mailto:i.gadzhiev.m@gmail.com}{i.gadzhiev.m@gmail.com} \\
\hline
\end{tabularx}
\caption{Code metadata.}
\end{table}
\section{Motivation and significance}
A researcher inspecting a scientific abstract may need to identify relevant
biological reactions in a curated catalog. Reaction Lens Systems supplies ranked proposals
for human review. Its engineering problem is maintaining an
inspectable relationship between input text, catalog entries and predictions.
Changing the order of candidates or the scoring block size should not silently
change the score assigned to an unchanged article--reaction pair.

We study three engineering questions: which scores should remain stable when the
catalog changes, how candidate blocking affects execution, and which component
dominates local prediction time. The contributions are (1) an explicit execution
contract derived from the global/local scorer and checked through candidate
transformations, (2) measurements of exact blocked scoring and an unblocked
reference, and (3) a reproducibility package linking aggregate results to recorded
observations and a separate installation check. Together, these distinguish
changes in rank caused by new competitors from changes in pairwise scores and
locate the main cost of complete-text prediction.
\section{Software description}
\subsection{Software architecture}
Reactome supplies the biological reaction catalog \citep{reactome2026}.
SPECTER and SPECTER2 provide citation-informed scientific document representations
\citep{cohan2020specter,singh2023scirepeval}. Reaction Lens Systems freezes its encoder and
uses a learned global/local head. We evaluate a fixed checkpoint to study
execution equivalence, component costs and inference reproducibility.

The processing path is English title/abstract, deterministic sentence spans,
frozen text encoding, projected reaction vectors, exact blocked scoring, then
ranked and threshold-selected output. SPECTER2 with its proximity adapter produces
768-dimensional features. Separate bias-free projections map article and reaction
features to 256 dimensions; the article projection is shared by global and
sentence features. Reaction vectors are prepared once during initialization.
The Python API, CLI and local HTTP service share the prediction implementation.
Measurements use sequential local calls. Loading verifies file hashes and array
shapes against the experiment manifest.

Metamorphic testing checks relations between transformed inputs and outputs when
individual expected outputs are difficult to specify \citep{segura2016}.
We apply this approach to candidate permutations, subsets and additions. The
scoring equation determines the expected relations; the experiments test whether
they hold in floating-point execution and across block configurations.
\subsection{Scoring and execution contracts}
Let $g$ and $s_j$ be the normalized projections of the global article and
valid sentence features through the shared learned article projection, and $r$
the normalized output of the separate learned reaction projection. The fixed head computes
\begin{equation}
 z(x,r)=a\left[g^\top r+w\max\left(0,\max_j s_j^\top r-n\right)\right]+b,
 \quad \widehat{y}(x,r)=\mathbf{1}[z(x,r)\geq\tau].
\end{equation}
Here the learned scale $a$, local weight $w$, null $n$, bias $b$, and development-fitted
threshold $\tau$ are held fixed. The implementation uses
$a=\exp(\min(\ell,\log 100))$, $w=\operatorname{softplus}(u)$ and
$n=\tanh(v)$ for learned scalars $\ell,u,v$. Padding is masked before the
sentence maximum. No reduction is taken over competing reactions.
Consequently, in real arithmetic, permuting candidates permutes outputs; removing
candidates preserves retained scores; and appending candidates preserves previous
scores. Fixed-threshold decisions inherit this property. Rankings and top-$k$
membership do not: new candidates can outrank existing ones. Updating descriptions,
model weights, preprocessing or the threshold also falls outside this contract.

For batch size $B$, sentence count $S$, candidate count $R$ and projected width
$P$, local similarities require $O(BSRP)$ arithmetic. Candidate blocks of width
$K$ bound the principal local-similarity intermediate to $O(BSK)$ instead of
$O(BSR)$. The final score matrix still uses $O(BR)$ storage, and the prepared
reaction bank uses $O(RP)$ storage. Blocking organizes exact scoring into smaller intermediate tensors while
preserving full catalog coverage.
\subsection{Software functionalities}
The local web interface accepts scientific text and displays ranked reactions and
their scores. The CLI exports JSON or CSV predictions, and the Python API and HTTP
service use the same prediction path. Model downloads are explicit and checked
against pinned hashes. The verification commands distinguish reconstruction of
recorded tables, cached-feature inference and rerunning complete-text measurements.
This separation lets a reader inspect the numerical evidence without downloading
models and then repeat the inference or measurement workload when desired.

\section{Illustrative examples}
\subsection{Local retrieval workflow}
From a checkout with Python 3.11 or 3.12 and uv installed, a user runs
\texttt{make setup} followed by \texttt{make serve}. The browser interface is
available at \url{http://127.0.0.1:8767/}. One demonstration input is
``STAT1 induces SMAD7 expression and inhibits TGF-beta signaling.'' The software
scores the complete catalog and displays proposals for inspection. This example
illustrates the retrieval workflow from input text to ranked catalog entries.
The reviewer can inspect API documentation at \url{http://127.0.0.1:8767/docs}.

Saving this demonstration sentence to a UTF-8 file named \texttt{abstract.txt},
the user can export the complete prediction payload with an empty title:
\begin{verbatim}
reaction-lens predict \
  --bundle artifacts/structured-seed23 \
  --encoder-dir artifacts/encoder \
  --title "" \
  --abstract-file abstract.txt \
  --format json --out predictions.json
\end{verbatim}
The input file must be supplied by the user, and the output path must not already
exist. The released README describes artifact preparation. Example results
are proposals, and the scores are not calibrated biological probabilities.

In a single illustrative CPU run using the released structured seed23 epoch105
bundle, the software scored all 15,692 reactions. The first three returned
reaction identifiers and raw logits were:
\begin{center}
\begin{tabular}{rrr}
\toprule
Rank & Reaction identifier & Raw logit \\
\midrule
1 & \texttt{R-HSA-9730967} & 7.286 \\
2 & \texttt{R-HSA-9736970} & 5.741 \\
3 & \texttt{R-HSA-2106591} & 5.580 \\
\bottomrule
\end{tabular}
\end{center}
This output demonstrates full-catalog execution and the format of ranked results.

\subsection{Reproducing the engineering evaluation}
\label{sec:protocol}
An AI assistant helped prepare measurement scripts and analyze results.
I am responsible for the content.

The recorded evidence can be checked with \texttt{make verify}, without model
downloads or third-party Python packages. After installation, \texttt{make
check-model} repeats the cached-feature evaluation and \texttt{make benchmark}
repeats the execution measurements. New observations are saved separately from
the published record. The evaluation protocol is as follows.

Experiments use the fixed structured-description checkpoint from seed 23 at epoch
105 and its checksummed cached development features. Model parameters and the
threshold are held fixed throughout the study.
All \CatalogCount{} catalog reactions are scored. We sample \SampleCount{} publications
without replacement using NumPy seed 20261005 and run CPU FP32 with two PyTorch
threads. The environment is macOS 26.6.1 on arm64, Python 3.12.14, PyTorch 2.14.1
and NumPy 2.5.3. We recorded the arm64 architecture in the benchmark metadata.

Block widths are 128, 1024 and the full catalog (15692).
Each block configuration receives five warmup queries and \RepeatCount{} timed repetitions
per sampled publication. Timings exclude encoding, bundle loading, initial
reaction projection and result serialization. Configurations run sequentially in
fixed order. Medians and p95 values summarize the recorded repeated measurements.

On the first repetition we check candidate permutation, a half-catalog subset,
appending 128 duplicate reaction vectors, and an unblocked reference. The
predeclared absolute logit tolerance is $10^{-5}$. The addition scenario reuses existing vectors, isolating candidate-axis
execution from text encoding. Threshold checks use $\tau=5.516382694244385$
without recalibration.

Memory is measured separately in a fresh process for each configuration using
the same padded batch of \MemoryBatch{} sampled publications, with maximum sentence count \MemorySpans{}.
We record process high-water resident memory including imports and loaded feature
arrays. Each configuration has one memory trial. Processes used for this comparison do not execute the
unblocked reference or the candidate-transformation probes.

Complete prediction measurements use \EndExamples{} English title/abstract
examples (PMIDs 10022914, 10047984, 10049780, 10067896 and 10069812) and three
repetitions after one warmup. Prediction includes text
encoding, scoring, ranking and payload construction. JSON serialization is recorded
separately; HTTP/network costs are excluded. Model setup is timed separately.

\section{Impact}
The software makes a complete-catalog reaction retrieval implementation available
for inspection and reuse, with an explicit distinction between changes in pairwise
scores and changes in rank. It provides a starting point for investigating how
catalog updates, candidate blocking and text representation affect scientific
retrieval execution. These capabilities support reproducible comparisons of
execution strategies and checks of implementation conformance to the specified
scoring behavior.

The reproducibility package supports inspection of the implementation,
reconstruction of the recorded tables and repetition of released-checkpoint
inference. The measurements inform execution choices by separating text-encoding
cost from full-catalog scoring and exposing the memory effect of candidate blocks.
\subsection{Candidate and execution equivalence}
The checks produce \ScoreComparisons{} scalar comparisons across
\ContractChecks{} combinations of queries, configurations and scenarios. The maximum observed
absolute logit error is \MaxError{}, with \DecisionFlips{} fixed-threshold flips.
All observed errors satisfy the predeclared tolerance. The checks cover 64
publications, three block configurations and four scenarios, confirming the
specified score and decision relations on the sampled inputs.

\subsection{Blocking and local execution cost}
\begin{table}[ht]
\centering
\begin{tabular}{rrrr}
\toprule
Block width & Median (ms) & p95 (ms) & Peak RSS (MiB)\\
\midrule
\SmallBlock & \SmallMedian & \SmallPninetyfive & \SmallMemory\\
\DefaultBlock & \DefaultMedian & \DefaultPninetyfive & \DefaultMemory\\
\FullBlock & \FullMedian & \FullPninetyfive & \FullMemory\\
\bottomrule
\end{tabular}
\caption{\raggedright Cached-feature single-publication scoring latency and separate padded-batch
process memory. These are different workloads. Every timing row scores the full
catalog. Latency covers 64 publications with three repetitions per
configuration; memory uses one fresh-process trial per configuration.}
\end{table}
The smallest block incurs greater observed latency. The default and unblocked
configurations have similar observed medians. The blocked configurations have lower observed
process peak RSS in the separate batch probe, consistent with the intermediate
storage bound.

The \EndRequests{} complete local predictions have median \EndMedian{} ms and p95
\EndPninetyfive{} ms. Encoding accounts for \EncodingShare\% of aggregate measured
prediction time. Setup takes \SetupSeconds{} s, and process peak RSS reaches
\EndMemory{} MiB. For these examples, improving only head execution offers limited benefit
to total latency because encoding dominates. The component breakdown identifies encoding as the primary target for
further latency experiments.

\subsection{Inference consistency in a separate installation}
We recomputed predictions in an isolated dependency environment using the
downloaded checkpoint and cached features for \PublicationCount{} development publications,
\KnownLinks{} known links and the full catalog. MAP is \MeasuredMAP{}, Hit@1 is
\MeasuredHit{}, and micro Recall@100 is \MeasuredRecall{}. Maximum absolute
difference from the stored release metrics is \MetricError{}.

For publication $i$, average precision sums the precision at each relevant
rank over the full catalog and divides by the number of distinct recorded
reaction labels for that publication, including labels absent from the catalog.
MAP averages this value over publications. Hit@1 is the fraction of publications
whose first result is recorded as relevant. Micro Recall@100 divides the total
recorded links recovered in the first 100 results by all recorded links across
publications. Equal logits retain catalog order. The absolute acceptance tolerance for the metric comparison is
$5\times10^{-5}$.

The check recovers the recorded metrics from the fixed checkpoint and cached
features. The complete-text latency experiment additionally exercises the local
encoder. These are separate checks of cached-feature inference and complete-text
execution.
\subsection{Interpretation of results}

The evaluation connects pairwise score stability, rank changes under catalog
updates and execution cost. Candidate independence follows from the pointwise
scoring equation; the experiments verify implementation conformance and quantify
execution costs. The latency and memory measurements identify distinct
optimization targets for text encoding and full-catalog scoring.

The reported timings characterize the workloads and environment specified in
Section~\ref{sec:protocol}. Configurations were executed in fixed order;
thermal and background-load effects were not controlled. RSS values are
whole-process measurements. The unblocked reference uses the same scoring
implementation, making the contract checks complementary to code inspection.
The addition scenario isolates candidate-axis execution by reusing reaction
vectors. The observed equivalence applies to the tested execution configurations.

Ranking metrics quantify agreement with recorded Reactome links. These labels
are incomplete and were used during model development; evaluation on unseen
publications is outside this execution study. The separate installation check
verifies released-checkpoint inference using cached features, while the
complete-text experiment exercises the local encoder.

\section{Conclusions}
Reaction Lens Systems connects a pointwise global/local scoring architecture to
explicit execution contracts. On the measured inputs, candidate transformations
and blocking preserve scores and fixed-threshold decisions. Blocking lowers
observed process peak memory in the batch probe, while text encoding dominates
complete local prediction time. The results provide a concrete procedure for
checking candidate-axis behavior and separating the costs of exact catalog
scoring from those of text representation.
\section*{Research materials}
The implementation, experiment protocol and recorded measurements are available
as \href{https://github.com/imrangadzhiev/reaction-lens-systems/tree/2140cc40ba3dc2af9ad0600cd2cc1624a47a2bf8}{research materials at revision \texttt{2140cc40ba3d}}.
Results in this article refer to that revision and the model and feature versions
identified by its manifests. 

The materials include \evidence{results/table3_contracts.csv}{contract checks},
\evidence{results/table1_scoring.csv}{block measurements},
\evidence{results/table2_end_to_end.csv}{component timings},
\evidence{results/fresh-evaluation.json}{the installation check}, and
\evidence{src/reaction_lens/evaluation.py}{metric definitions in code}.
The \evidence{results/system-benchmark.json}{raw experiment record} and
\evidence{NUMBERS.csv}{numerical index} connect aggregate values to observations.
The protocol is described in Section~\ref{sec:protocol}; the
\evidence{PROTOCOL.md}{accompanying protocol} and
\evidence{README.md}{execution instructions} provide operational details for
reconstructing tables, rerunning inference and repeating measurements. These are
separate operations. Source checksums identify the measured scorer, encoder,
pipeline and benchmark. Code is licensed under MIT and the manuscript under
CC BY 4.0.

\par\addvspace{1.5\baselineskip}\noindent\begin{minipage}{\textwidth}
\section*{Author declarations}
\textbf{Funding:} No external funding was received for this work.

\textbf{Declaration of competing interests:} I declare no competing financial
interests or personal relationships that could have influenced this work.

\textbf{CRediT author contribution statement:} Imran Gadzhiev: Conceptualization,
Methodology, Software, Investigation, Validation, Writing -- original draft,
Writing -- review \& editing.

\section*{Declaration of generative AI and AI-assisted technologies in the manuscript preparation process}
I used Codex (OpenAI) to assist with manuscript preparation under my direction.
I take full responsibility for the content.
\end{minipage}

\begingroup\small
\bibliographystyle{elsarticle-num}
\bibliography{references}
\endgroup
\end{document}
